\documentclass[aspectratio=169,11pt]{beamer}
\usetheme{Madrid}
\usecolortheme{default}
\setbeamertemplate{navigation symbols}{}
\usepackage[utf8]{inputenc}
\usepackage[T1]{fontenc}
\usepackage{amsmath,amssymb,booktabs,array,graphicx}

\newcommand{\ind}{\mathbf{1}}
% Placeholder box for figures: replace with \includegraphics when adding the plot
\newcommand{\figplaceholder}[1]{%
  \fbox{\parbox[c][0.5\textheight][c]{0.9\linewidth}{\centering\small #1}}}

\title[Temperature and electricity demand]{Temperature and Electricity Demand in Spain}
\subtitle{Case Study II: Heat waves and the electricity market, step 2}
\author{}
\date{October 2026}

\begin{document}

\begin{frame}
  \titlepage
\end{frame}

% ===========================================================================
\section{Question and data}
% ===========================================================================

\begin{frame}{Question}
  \begin{itemize}
    \item Electricity demand has strong regular patterns: hour of day, day of week,
          holidays, season.
    \item \textbf{Question:} once those patterns are removed, how much does a
          temperature shock move demand?
    \item Two demand series, always estimated at hourly frequency:
      \begin{itemize}
        \item \textbf{Aggregate (peninsula):} all consumers. System-level response,
              the input for the market simulation (counterfactual prices).
        \item \textbf{Small consumers (47 provinces):} mostly households. Matched to
              local weather; allows heterogeneous effects across provinces.
      \end{itemize}
  \end{itemize}
\end{frame}

\begin{frame}{Data: demand}
  \small
  \begin{tabular}{@{}l>{\raggedright\arraybackslash}p{5.6cm}>{\raggedright\arraybackslash}p{5.6cm}@{}}
    \toprule
    & Aggregate (peninsula) & Small consumers (province) \\
    \midrule
    Variable     & $\ln D_t$: log hourly demand & $\ln D_{i,t}$: log hourly demand of province $i$ \\
    Unit         & MW (hourly mean) & kWh per hour \\
    Spatial unit & Peninsula, one series & 47 peninsular provinces \\
    Mean         & 26{,}841 MW & 165{,}150 kWh per province-hour (7{,}762 MW summed) \\
    Coverage     & All consumers & Small consumers, about 29\% of demand \\
    Observations & 40{,}148 hours & 1{,}886{,}956 province-hours \\
    \bottomrule
  \end{tabular}
  \vspace{4pt}
  \begin{itemize}\footnotesize
    \item Community-level series are not used: they equal peninsular demand times a
          share fixed within each month, so they carry no independent hourly variation.
  \end{itemize}
\end{frame}

\begin{frame}{Data: temperature, calendar and sample}
  \begin{columns}[T]
    \begin{column}{0.52\textwidth}
      \textbf{Degree-days} (daily, by province, $^{\circ}$C$\cdot$day)
      \[
        HDD_{i,d} = \max(14^{\circ}\mathrm{C} - T_{i,d},\,0), \quad
        CDD_{i,d} = \max(T_{i,d} - 21^{\circ}\mathrm{C},\,0)
      \]
      \vspace{-6pt}
      \begin{itemize}\small
        \item $T_{i,d}$: daily mean temperature of province $i$ on day $d$.
        \item Same value for all 24 hours of a day.
        \item Peninsula: weighted mean of the 47 provinces (weights = demand shares).
      \end{itemize}
    \end{column}
    \begin{column}{0.46\textwidth}
      \textbf{Calendar}
      \begin{itemize}\small
        \item Hour of day, day type (Mon--Sun or holiday), month, year.
        \item Holidays: 9 fixed national days, Good Friday, Holy Thursday (50 days).
      \end{itemize}
      \textbf{Sample}
      \begin{itemize}\small
        \item 1 June 2021 -- 31 December 2025: after COVID and the June 2021 2.0TD
              time-of-use tariff.
        \item 28--29 April 2025 (blackout) excluded.
      \end{itemize}
    \end{column}
  \end{columns}
\end{frame}

% ===========================================================================
\section{Baseline regressions}
% ===========================================================================

\begin{frame}{Notation}
  \begin{columns}[T]
    \begin{column}{0.48\textwidth}
      \begin{itemize}\small
        \item $t$: an hourly observation (local time)
        \item $i$: province, $i = 1,\dots,47$
        \item $d(t)$: calendar day of hour $t$
        \item $h(t) \in \{0,\dots,23\}$: hour of day
        \item $m(t) \in \{1,\dots,12\}$: month of year
      \end{itemize}
    \end{column}
    \begin{column}{0.48\textwidth}
      \begin{itemize}\small
        \item $\tau(t)$: day type (Mon, \dots, Sun, holiday)
        \item $y(t)$: year; $ym(t)$: year-month
        \item $\alpha, \gamma, \delta$: fixed effects (one intercept per cell)
        \item $\ind\{h(t)=k\}$: equals 1 if $t$ is in hour $k$, 0 otherwise
        \item $\beta_H, \beta_C$: effects of HDD and CDD
      \end{itemize}
    \end{column}
  \end{columns}
  \vspace{8pt}
  \small Coefficients $\times 100$ $\approx$ percent change in demand per degree-day.
\end{frame}

\begin{frame}{Baseline specification (M2)}
  \textbf{Peninsula}
  \[
    \ln D_t = \beta_H\, HDD_{d(t)} + \beta_C\, CDD_{d(t)}
      + \alpha_{h(t),m(t)} + \gamma_{h(t),\tau(t)} + \delta_{y(t)} + \varepsilon_t
  \]
  \textbf{Small consumers}
  \[
    \ln D_{i,t} = \beta_H\, HDD_{i,d(t)} + \beta_C\, CDD_{i,d(t)}
      + \alpha_{i,h(t),m(t)} + \gamma_{i,h(t),\tau(t)} + \delta_{i,y(t)} + \varepsilon_{i,t}
  \]
  \begin{itemize}\small
    \item $\alpha$: hour $\times$ month: normal daily load shape in each month.
    \item $\gamma$: hour $\times$ day type: weekly cycle and holiday shape.
    \item $\delta$: year: growth, rooftop solar, 2022 price crisis.
    \item Identification: days hotter or colder than usual for that province and month.
    \item SE clustered by week (peninsula); by province and week (provinces).
  \end{itemize}
\end{frame}

\begin{frame}{Specification ladder}
  \small
  \begin{tabular}{@{}lll@{}}
    \toprule
    Model & Fixed effects (peninsula; provinces add $\times$ province) & Idea \\
    \midrule
    M1  & $\alpha_h + \alpha_\tau + \alpha_m + \alpha_y$ & Additive calendar controls \\
    M2  & $\alpha_{h,m} + \gamma_{h,\tau} + \delta_y$ & Baseline \\
    M3  & $\alpha_{h,m} + \gamma_{h,\tau} + \delta_{ym}$ & Within-month shocks only \\
    M2w & M2, provinces weighted by mean demand & Effect on total small-consumer demand \\
    M4  & M2, with HDD and CDD effects by hour of day & Hourly response profile \\
    \bottomrule
  \end{tabular}
  \vspace{8pt}
  \textbf{M4 (peninsula; provinces analogous with index $i$)}
  \[
    \ln D_t = \sum_{k=0}^{23} \beta^{C}_{k}\, \ind\{h(t)=k\}\, CDD_{d(t)}
      + \sum_{k=0}^{23} \beta^{H}_{k}\, \ind\{h(t)=k\}\, HDD_{d(t)}
      + \alpha_{h(t),m(t)} + \gamma_{h(t),\tau(t)} + \delta_{y(t)} + \varepsilon_t
  \]
  \footnotesize $\beta^{C}_{k}$: percent change in demand at hour $k$ per extra CDD that day.
\end{frame}

\begin{frame}{Baseline results}
  \centering\small
  \begin{tabular}{@{}lccccc@{}}
    \toprule
    & M1 & M2 & M3 & M2w & M4 \\
    \midrule
    \multicolumn{6}{@{}l}{\textit{Peninsula (aggregate)}} \\
    \quad CDD21 & 2.81*** & 2.81*** & 2.75*** & -- & by hour \\
                & (0.11)  & (0.11)  & (0.11)  &    &         \\
    \quad HDD14 & 1.77*** & 1.77*** & 1.60*** & -- & by hour \\
                & (0.17)  & (0.17)  & (0.17)  &    &         \\
    \quad $R^2$ & 0.850 & 0.927 & 0.934 & -- & 0.930 \\[3pt]
    \multicolumn{6}{@{}l}{\textit{Small consumers (province panel)}} \\
    \quad CDD21 & 3.58*** & 3.35*** & 3.25*** & 4.06*** & by hour \\
                & (0.23)  & (0.22)  & (0.22)  & (0.22)  &         \\
    \quad HDD14 & 2.08*** & 2.24*** & 2.08*** & 2.78*** & by hour \\
                & (0.19)  & (0.15)  & (0.14)  & (0.20)  &         \\
    \quad $R^2$ & 0.984 & 0.995 & 0.996 & 0.995 & 0.995 \\
    \bottomrule
  \end{tabular}
  \vspace{4pt}

  \footnotesize\raggedright
  Coefficients $\times 100$ (\% per degree-day). SE in parentheses: clustered by week
  (peninsula) or by province and week (provinces). *** $p<0.01$.
  N = 40{,}148 hours and 1{,}886{,}956 province-hours. M1--M3 for the
  peninsula use small-consumer weights for national degree-days.
\end{frame}

\begin{frame}{Interpretation}
  \begin{columns}[T]
    \begin{column}{0.48\textwidth}
      \textbf{Aggregate demand (M2)}
      \begin{itemize}\small
        \item +1 CDD: +2.8\%, about +754 MW.
        \item +1 HDD: +1.8\%, about +475 MW.
        \item Very hot day (95th pct., CDD = 5.9):
              $e^{0.0281 \times 5.92} - 1 \approx +18\%$, close to +5 GW.
      \end{itemize}
    \end{column}
    \begin{column}{0.48\textwidth}
      \textbf{Small consumers (M2w)}
      \begin{itemize}\small
        \item +1 CDD: +4.1\%, about +315 MW.
        \item +1 HDD: +2.8\%, about +216 MW.
        \item 29\% of demand, but about 42\% of the system's response to heat.
      \end{itemize}
    \end{column}
  \end{columns}
  \vspace{8pt}
  \begin{itemize}\small
    \item M1 and M2 give the same peninsula coefficient: CDD is constant within a day,
          so hour interactions do not change its identifying variation; they raise $R^2$.
    \item For provinces, M1 and M2 differ because M2 interacts the fixed effects with province.
  \end{itemize}
\end{frame}

\begin{frame}{Hourly profile (M4)}
  \centering
  \figplaceholder{Insert \texttt{output/baseline/hourly\_effects\_M4.png}}
  % \includegraphics[width=0.95\linewidth]{hourly_effects_M4.png}
  \vspace{4pt}
  \begin{itemize}\small
    \item Cooling effect peaks at 16--17h: 4.1\% (peninsula), 5.2\% (small consumers).
    \item Small consumers also respond at night (3.3--3.6\% at 0--4h vs.\ 2.2--2.8\%).
  \end{itemize}
\end{frame}

% ===========================================================================
\section{Climate zones}
% ===========================================================================

\begin{frame}{Using province-level weather in the peninsula regression}
  \textbf{Why not all 47 provinces' HDD and CDD as regressors?}
  \begin{itemize}\small
    \item 94 temperature coefficients (2{,}256 by hour), identified from only 1{,}673
          days of weather.
    \item Strong collinearity: within a month, provincial CDD correlate 0.58 on average.
    \item Individual coefficients would be noisy and uninterpretable.
  \end{itemize}
  \textbf{But regional variation exists}
  \begin{itemize}\small
    \item One common factor explains 60\% of within-month CDD variation; a second,
          regional factor a further 13\%.
  \end{itemize}
  \textbf{Solution:} group provinces into four climate zones, each with its own coefficient.
  The data then choose the weights.
\end{frame}

\begin{frame}{Climate zones: provinces and demand shares}
  \scriptsize
  \begin{tabular}{@{}l>{\raggedright\arraybackslash}p{9.9cm}r@{}}
    \toprule
    Zone & Communities (share of peninsular demand) and provinces & Share \\
    \midrule
    North & Galicia (6.0\%): La Coruña, Lugo, Orense, Pontevedra;
            Asturias (3.6\%); Cantabria (1.6\%);
            País Vasco (6.3\%): Álava, Guipúzcoa, Vizcaya & 17.5\% \\[2pt]
    Interior & Castilla y León (5.6\%): Ávila, Burgos, León, Palencia, Salamanca,
               Segovia, Soria, Valladolid, Zamora; La Rioja (0.7\%); Navarra (2.0\%);
               Aragón (4.3\%): Huesca, Teruel, Zaragoza; Madrid (11.7\%);
               Castilla-La Mancha (5.1\%): Albacete, Ciudad Real, Cuenca, Guadalajara, Toledo;
               Extremadura (2.1\%): Badajoz, Cáceres & 31.5\% \\[2pt]
    Mediterranean & Cataluña (18.9\%): Barcelona, Gerona, Lérida, Tarragona;
                    Comunidad Valenciana (11.5\%): Alicante, Castellón, Valencia;
                    Murcia (3.9\%) & 34.3\% \\[2pt]
    South & Andalucía (16.7\%): Almería, Cádiz, Córdoba, Granada, Huelva, Jaén,
            Málaga, Sevilla & 16.7\% \\
    \bottomrule
  \end{tabular}
  \vspace{4pt}

  \raggedright
  Shares: community share of peninsular demand, June 2021 -- December 2025.
  Check: clustering provinces on daily temperature deviations gives similar groups
  for the interior; results are robust to using the data-driven clusters.
\end{frame}

\begin{frame}{Zone specification}
  \textbf{Weights} (total demand):
  $w_p = s_{c(p)} \times \bar{q}_p / \sum_{p' \in c(p)} \bar{q}_{p'}$,
  with $s_c$ the community's demand share and $\bar{q}_p$ mean small-consumer demand.
  \[
    CDD_{z,d} = \frac{\sum_{p \in z} w_p\, CDD_{p,d}}{\sum_{p \in z} w_p}
    \quad\text{(same for HDD)}
  \]
  \textbf{Zone model by hour (Z4)}
  \[
    \ln D_t = \sum_{z}\sum_{k=0}^{23} \beta^{C}_{z,k}\, \ind\{h(t)=k\}\, CDD_{z,d(t)}
      + \sum_{z}\sum_{k=0}^{23} \beta^{H}_{z,k}\, \ind\{h(t)=k\}\, HDD_{z,d(t)}
      + \alpha_{h(t),m(t)} + \gamma_{h(t),\tau(t)} + \delta_{y(t)} + \varepsilon_t
  \]
  \begin{itemize}\small
    \item $z \in \{$North, Interior, Mediterranean, South$\}$; $4 \times 24 \times 2 = 192$
          temperature coefficients.
    \item $\beta^{C}_{z,k}$: percent change in national demand at hour $k$ per extra CDD
          in zone $z$, other zones held fixed.
    \item Z2: one coefficient per zone. Test (T2): national degree-days plus zone
          deviations; Wald test that deviations are zero.
  \end{itemize}
\end{frame}

\begin{frame}{Zone results (pooled over hours)}
  \centering\scriptsize
  \begin{tabular}{@{}lccc@{}}
    \toprule
    & (1) National (N2) & (2) Zones (Z2) & (3) Test (T2) \\
    \midrule
    \multicolumn{4}{@{}l}{\textit{CDD21}} \\
    \quad National      & 2.896*** (0.118) &                    & 2.600*** (0.182) \\
    \quad North         &                  & $-$0.022 (0.204)   &                  \\
    \quad Interior      &                  & 0.649*** (0.124)   & 0.689 (0.446)    \\
    \quad Mediterranean &                  & 1.457*** (0.116)   & 1.500*** (0.420) \\
    \quad South         &                  & 0.515*** (0.113)   & 0.536*** (0.207) \\
    \multicolumn{4}{@{}l}{\textit{HDD14}} \\
    \quad National      & 1.712*** (0.164) &                    & 2.095*** (0.250) \\
    \quad North         &                  & 0.213 (0.228)      &                  \\
    \quad Interior      &                  & 0.015 (0.303)      & $-$0.368 (0.688) \\
    \quad Mediterranean &                  & 1.345*** (0.248)   & 0.928** (0.451)  \\
    \quad South         &                  & 0.523* (0.271)     & 0.320 (0.274)    \\
    \midrule
    Wald $F(6,239)$, deviations $= 0$ & & & 8.54 ($p<0.001$) \\
    $R^2$               & 0.927 & 0.929 & 0.929 \\
    \bottomrule
  \end{tabular}
  \vspace{3pt}

  \raggedright\tiny
  Coefficients $\times 100$. Hourly log peninsular demand, N = 40{,}148. FE: hour $\times$
  month, hour $\times$ day type, year. SE clustered by week. In (3), zone rows are deviations
  from national degree-days (North omitted). * $p<0.10$, ** $p<0.05$, *** $p<0.01$.
\end{frame}

\begin{frame}{Relative importance: demand share vs.\ estimated effect}
  \centering\small
  \begin{tabular}{@{}lcccc@{}}
    \toprule
    Zone & Demand share & CDD coef. & Share of zone CDD effect & Share of zone HDD effect \\
    \midrule
    Mediterranean & 34.3\% & 1.46 & 56\% & 64\% \\
    Interior      & 31.5\% & 0.65 & 25\% & 1\%  \\
    South         & 16.7\% & 0.52 & 20\% & 25\% \\
    North         & 17.5\% & $-$0.02 & $-$1\% & 10\% \\
    \midrule
    Sum           & 100\%  & 2.60 & 100\% & 100\% \\
    \bottomrule
  \end{tabular}
  \vspace{6pt}
  \begin{itemize}\small
    \item Share of effect = zone coefficient / sum of zone coefficients (Z2).
    \item Mediterranean heat matters far more than its demand share; northern heat
          has no measurable effect.
    \item Fixed national weights rejected, but a uniform +1 CDD gives a similar
          total: 2.60\% (zones) vs.\ 2.90\% (national).
  \end{itemize}
\end{frame}

\begin{frame}{Zone effects by hour (Z4)}
  \centering\scriptsize
  \begin{tabular}{@{}lcccccc@{}}
    \toprule
    CDD, \% per zone degree-day & 0h & 4h & 8h & 12h & 16h & 20h \\
    \midrule
    North         & $-$0.01 (0.19) & 0.08 (0.20) & 0.01 (0.30) & 0.02 (0.34) & $-$0.05 (0.31) & $-$0.17 (0.31) \\
    Interior      & 0.26 (0.13) & 0.15 (0.12) & 0.32 (0.19) & 0.82 (0.22) & 1.20 (0.21) & 1.10 (0.19) \\
    Mediterranean & 1.65 (0.14) & 1.34 (0.13) & 1.08 (0.15) & 1.71 (0.23) & 1.91 (0.20) & 1.05 (0.20) \\
    South         & 0.75 (0.12) & 0.65 (0.12) & 0.14 (0.16) & 0.13 (0.19) & 0.66 (0.18) & 0.68 (0.16) \\
    \midrule
    \multicolumn{7}{@{}l}{$R^2$: Z4 0.932; national by hour (N4) 0.930} \\
    \bottomrule
  \end{tabular}
  \vspace{6pt}
  \begin{itemize}\small
    \item Mediterranean: high at all hours, including night (humid coast).
    \item Interior: afternoon peak (15--17h). South: later, around 18h.
    \item Wald test of fixed weights by hour: $F(144,239) = 8.4$, $p<0.001$.
  \end{itemize}
\end{frame}

\begin{frame}{Zone effects by hour: plot}
  \centering
  \figplaceholder{Insert \texttt{output/zones/hourly\_zone\_effects\_Z4.png}}
  % \includegraphics[width=0.95\linewidth]{hourly_zone_effects_Z4.png}
  \vspace{4pt}

  \small Percent change in national demand per zone degree-day, by hour; 95\% bands.
\end{frame}

% ===========================================================================
\section{Summary}
% ===========================================================================

\begin{frame}{Summary and next steps}
  \begin{columns}[T]
    \begin{column}{0.48\textwidth}
      \textbf{Findings}
      \begin{itemize}\small
        \item Heat raises demand more than cold: +2.8\% vs.\ +1.8\% per degree-day.
        \item Small consumers respond more (+4.1\% per CDD) and drive about 42\% of the
              heat response.
        \item Response peaks at 16--18h.
        \item Mediterranean heat is the main driver of national demand.
      \end{itemize}
    \end{column}
    \begin{column}{0.48\textwidth}
      \textbf{Next steps}
      \begin{itemize}\small
        \item Temperature bins: nonlinear effects at extreme heat.
        \item Heat-wave indicator and lags.
        \item Province heterogeneity (province $\times$ CDD).
        \item Regional holidays; humidity, solar radiation, wind.
      \end{itemize}
    \end{column}
  \end{columns}
\end{frame}

\end{document}
